\input{ieee-preamble}
\hypersetup{
  pdftitle={Entropy-Minimizing Optimal Control of a Helicopter},
  pdfauthor={Santiago Restrepo},
  pdfsubject={Proposed thermodynamic optimal-control method; one-page technical summary},
  pdfkeywords={helicopter, entropy generation, exergy, nonlinear MPC, optimal control}
}

\title{Entropy-Minimizing Control of a Helicopter}
\author{\IEEEauthorblockN{Santiago Restrepo}
\IEEEauthorblockA{\href{mailto:contact@waajacu.com}{contact@waajacu.com}\quad
\url{https://waajacu.com/}\\[3pt]
\textit{Proposed thermodynamic optimal-control method}}}
\date{}

\begin{document}
\maketitle

\begin{abstract}
This note proposes entropy-generation-minimizing control for a conventional
helicopter near hover. A nonlinear predictive controller seeks to reduce
irreversible losses while completing a prescribed flight task. The method
combines coupled rotorcraft dynamics, thermodynamic loss maps, and explicit
operating constraints; it requires vehicle-specific identification and
closed-loop validation.
\end{abstract}

\begin{IEEEkeywords}
Helicopter, entropy generation, exergy, nonlinear optimal control.
\end{IEEEkeywords}

\balance
\section{Define the Mission and Dynamics}
Specify a fixed-duration hover-to-hover maneuver, payload, flight corridor,
and endpoint tolerances. This prevents apparent savings obtained by doing
less work. Assume a single main rotor, a tail rotor, and governed rotor speed.
\begin{equation}
\begin{aligned}
\dot{x}&=f(x,u,w),\\
u&=[\theta_0,\theta_{1c},\theta_{1s},\theta_t]^{\mathsf T}.
\end{aligned}
\end{equation}
The inputs are main collective, two cyclic pitches, and tail collective.
State $x$ contains position, velocity, attitude, body rates, and relevant
rotor flapping, inflow, actuator, and thermal states; $w$ represents wind.
Identify the coupled force and moment model around trim. Include governor
dynamics when rotor-speed transients matter.

\section{Account for Physical Entropy}
Choose a boundary covering the aircraft, propulsion, and specified wake
relaxation. Define entropy generation rate $\sigma$ in W/K. At fixed ambient
temperature $T_0$, destroyed exergy equals $T_0$ times generated
entropy~\cite{philipp2024}. A component model gives
\begin{equation}
\begin{aligned}
\sigma&=\sum_j P_{\mathrm{irr},j}/T_j+\sigma_{\mathrm{other}}\geq 0,\\
P_{\mathrm{ex,dest}}&=T_0\sigma.
\end{aligned}
\end{equation}
Here $P_{\mathrm{irr},j}$ is locally dissipated power at absolute temperature
$T_j$; $\sigma_{\mathrm{other}}$ includes heat-transfer, mixing, and chemical
or electrochemical irreversibility as applicable. Fit these terms from
component balances and measured loss maps without double counting.

Rotor-induced power initially leaves as organized wake energy. Count it as
entropy generation only when wake relaxation is included; otherwise retain
it as outgoing exergy. Shaft power divided by ambient temperature is
therefore not a general entropy model. Track stored mechanical and thermal
energy so the horizon cannot hide losses.

\section{Connect Rotor Loads to the Cost}
A near-hover power estimate for each rotor is~\cite{johnson2015}
\begin{equation}
P_r\approx\kappa_r F_r^{3/2}/(2\rho A_r)^{1/2}+P_{\mathrm{profile},r}.
\end{equation}
$F_r$ is thrust, $A_r$ disk area, $\rho$ air density, and $\kappa_r$ the
induced-power correction. Sum main- and tail-rotor requirements and
drivetrain losses. Use calibrated inflow and drag maps for moving flight;
the hover relation does not cover stall or vortex-ring descent. These power
terms feed the entropy accounting in Section~II.

\section{Solve a Constrained Flight Problem}
Discretize the dynamics over $N$ steps of duration $\Delta t$ and minimize
total generated entropy:
\begin{equation}
\begin{aligned}
\min_U\ J&=\sum_{k=0}^{N-1}\Delta t\,\sigma(x_k,u_k,w_k),\\
x_{k+1}&=f_{\Delta t}(x_k,u_k,w_k),\\
e_k&\in E_k,\quad x_N\in X_f,\quad(x_k,u_k,\Delta u_k)\in C.
\end{aligned}
\end{equation}
$U$ is the input sequence, $x_0$ the estimated state, and $w_k$ the wind
forecast. $E_k$ bounds tracking error; $X_f$ specifies terminal flight and
mechanical/thermal states. $C$ limits pitch commands and rates, rotor speed,
shaft torque/power, temperatures, attitude, clearance, and the identified
flight envelope. Choose a terminal set compatible with a stabilizing
controller; entropy minimization alone does not ensure stability.

\section{Implement and Validate}
First identify the model and solve offline by direct multiple shooting or
collocation (e.g., CasADi~\cite{casadi}). Then estimate states, warm-start the
optimization, apply its first input, and repeat as nonlinear model
predictive control. Check feasibility and computation time; retain a
stabilizing fallback.

Compare the same maneuver with PID/LQR under wind, mass, and model errors.
Report tracking error, integrated entropy or a declared power proxy, input
energy, peak torque, and violations. Hover still requires positive power;
no efficiency gain is established without these tests.

\begingroup
\raggedright
\bibliographystyle{IEEEtran}
\bibliography{references}
\endgroup
\end{document}
